% Generated from locale/tr/content/set-theory/spine/recursion.tex; do not hand-edit.


\olfileid[tr]{sth}{spine}{recursion}
\olsection{Sonlu Ötesi Özyineleme Teoremleri}

Önce \olref[valpha]{defValphas} tanımının başarılı olduğunu
doğrulamalıyız. Daha genel olarak, sonlu ötesi özyinelemeyle bir tanım
vermeye yönelik her girişimin başarıya ulaşacağını ispatlamamız gerekir.
Bu kesimin amacı budur.

\emph{Uyarı: Bu malzeme zordur}. Bununla birlikte ana fikir oldukça
basittir: Sonlu Ötesi Tümevarım ile Yerine Koyma, sonlu ötesi özyinelemenin
(çeşitli sürümlerinin) meşruluğunu güvence altına alır.\footnote{Hatırlatma:
açıkça aksi belirtilmedikçe bütün formüller ve terimler parametreler
içerebilir.}

\begin{defn}
	$\tau(x)$ bir terim, $f$ bir fonksiyon ve $\alpha$ bir sıra sayısı olsun.
	Hem $\dom{f}=\alpha$ hem de
	$(\forall\beta\in\alpha)f(\beta)=
	\tau(\funrestrictionto{f}{\beta})$ ise $f$'ye $\tau$ için bir
	\emph{$\alpha$-yaklaştırması} deriz.
\end{defn}

\begin{lem}[Sınırlı Özyineleme]\ollabel{transrecursionfun}
Her $\tau(x)$ terimi ve her $\alpha$ sıra sayısı için $\tau$'nun biricik bir
$\alpha$-yaklaştırması vardır.
\end{lem}
\begin{proof}
Her $\gamma\leq\alpha$ için biricik bir $\gamma$-yaklaştırması bulunduğunu
göstereceğiz.

Önce biricikliği kuruyoruz. $g$ ve $h$ sırasıyla $\gamma$- ve
$\delta$-yaklaştırmaları olsun. Argümanları üzerinde sonlu ötesi tümevarım,
her $\beta\in\dom{g}\cap\dom{h}=\gamma\cap\delta=
\min(\gamma,\delta)$ için $g(\beta)=h(\beta)$ olduğunu gösterir. Dolayısıyla
yaklaştırmalarımız (varlarsa) biriciktir ve bütün ortak değerlerde uyuşur.

Varlığı kurmak için şimdi $\delta\leq\alpha$ sıra sayıları üzerinde basit
sonlu ötesi tümevarım (\olref[ordinals][opps]{simpletransrecursion})
kullanıyoruz.

Boş fonksiyon açıkça bir $\emptyset$-yaklaştırmasıdır.

$g$ bir $\gamma$-yaklaştırmasıysa
$g\cup\{\tuple{\gamma,\tau(g)}\}$ bir
$\ordsucc{\gamma}$-yaklaştırmasıdır.

$\gamma$ bir limit sıra sayısı ve her $\delta<\gamma$ için $g_\delta$ bir
$\delta$-yaklaştırmasıysa
$g=\bigcup_{\delta\in\gamma}g_\delta$ olsun. Çeşitli $g_\delta$'larımız
bütün değerlerde uyuştuğu için bu bir fonksiyondur. Ayrıca
$\delta\in\gamma$ ise
$g(\delta)=g_{\ordsucc{\delta}}(\delta)=
\tau(\funrestrictionto{g_{\ordsucc{\delta}}}{\delta})=
\tau(\funrestrictionto{g}{\delta})$ olur.

Bu, sonlu ötesi tümevarımla ispatı tamamlar.
\end{proof}

Bir fonksiyon yerine bir \emph{terim} tanımlamamıza izin verirsek önceki
sonuçtaki $\alpha$ sınırını kaldırabiliriz. Aşağıdaki sonucun ifade ve
ispatında $\sigma$ bir terim olduğunda
$\funrestrictionto{\sigma}{\alpha}=
\Setabs{\tuple{\beta,\sigma(\beta)}}{\beta\in\alpha}$ olsun.

\begin{thm}[Genel Özyineleme]\ollabel{transrecursionschema}
Her $\tau(x)$ terimi için, her $\alpha$ sıra sayısında
$\sigma(\alpha)=\tau(\funrestrictionto{\sigma}{\alpha})$ olacak biçimde
bir $\sigma(x)$ terimini açıkça tanımlayabiliriz.
\end{thm}

\begin{proof}
Her $\alpha$ için \olref{transrecursionfun} gereği $\tau$'nun biricik bir
$\alpha$-yaklaştırması $f_\alpha$ vardır. $\sigma(\alpha)$'yı
$f_{\ordsucc{\alpha}}(\alpha)$ olarak tanımlayın. Şimdi:
	\begin{align*}
		\sigma(\alpha) &=
		f_{\ordsucc{\alpha}}(\alpha) \\&=
		\tau(\funrestrictionto{f_{\ordsucc{\alpha}}}{\alpha}) \\&=
		\tau(\Setabs{\tuple{\beta, f_{\ordsucc{\alpha}}(\beta)}}{\beta \in \alpha}) \\&=
		\tau(\Setabs{\tuple{\beta, f_{\ordsucc{\beta}}(\beta)}}{\beta \in \alpha})\\&=
		\tau(\funrestrictionto{\sigma}{\alpha}).
	\end{align*}
Burada bütün $\beta<\alpha$ için
$f_{\ordsucc{\beta}}(\beta)=f_{\ordsucc{\alpha}}(\beta)$ olduğu,
\olref{transrecursionfun} sonucundaki gibi kullanılmıştır.
\end{proof}
\noindent
\olref{transrecursionschema} sonucunun bir \emph{şema} olduğuna dikkat
edin. Kritik olarak $\sigma$'nın bir fonksiyon, yani belirli türden bir
\emph{küme} tanımlamasını bekleyemeyiz; çünkü o zaman $\dom{\sigma}$ bütün
sıra sayılarının kümesi olur ve Burali--Forti Paradoksu'yla
(\olref[ordinals][basic]{buraliforti}) çelişirdi.

Yine de \olref{transrecursionschema} sonucunun $V_\alpha$ tanımımızı
doğruladığını göstermemiz gerekir. Bu hemen açık olmayabilir; fakat son ve
basit bir sonlu ötesi özyineleme sürümüyle görünür hâle gelecektir.

\begin{thm}[Basit Özyineleme]\ollabel{simplerecursionschema}
Her $\tau(x)$ ve $\theta(x)$ terimi ile her $A$ kümesi için şu koşulları
sağlayan bir $\sigma(x)$ terimini açıkça tanımlayabiliriz:
\begin{align*}
	\sigma(\emptyset) &= A\\
	\sigma(\ordsucc{\alpha}) &= \tau(\sigma(\alpha)) &&
		\text{her $\alpha$ sıra sayısı için}\\
	\sigma(\alpha) &= \theta(\ran{\funrestrictionto{\sigma}{\alpha}})&&
	\text{$\alpha$ bir limit sıra sayısı olduğunda.}
\end{align*}
\end{thm}

\begin{proof}
Önce $\xi(x)$ terimini şöyle tanımlıyoruz:
%t $\xi(x) = A$ if $x$ is not a function whose domain is an ordinal;
%otherwise let:
\[
	\xi(x) =
	\begin{cases}
			A &\text{$x$, tanım kümesi bir sıra sayısı olan}\\
			&\hspace{1em}\text{bir fonksiyon değilse; aksi hâlde:}\\
			\tau(x(\alpha)) & \text{$\dom{x} = \ordsucc{\alpha}$ ise}\\
			\theta(\ran{x}) & \text{$\dom{x}$ bir limit sıra sayısıysa.}
	\end{cases}
\]
\olref{transrecursionschema} gereği her $\alpha$ sıra sayısı için
$\sigma(\alpha)=\xi(\funrestrictionto{\sigma}{\alpha})$ olan bir
$\sigma(x)$ terimi vardır; ayrıca
$\funrestrictionto{\sigma}{\alpha}$, tanım kümesi $\alpha$ olan bir
fonksiyondur. Basit sonlu ötesi tümevarımla
(\olref[ordinals][opps]{simpletransrecursion}) $\sigma$'nın gereken
özelliklere sahip olduğunu gösteriyoruz.

Önce $\sigma(\emptyset)=\xi(\emptyset)=A$ olur.

Sonra $\sigma(\ordsucc{\alpha})=
\xi(\funrestrictionto{\sigma}{\ordsucc{\alpha}})=
\tau(\funrestrictionto{\sigma}{\ordsucc{\alpha}}(\alpha))=
\tau(\sigma(\alpha))$ olur.

Son olarak $\alpha$ bir limit olduğunda
$\sigma(\alpha)=\xi(\funrestrictionto{\sigma}{\alpha})=
\theta(\ran{\funrestrictionto{\sigma}{\alpha}})$ olur.
\end{proof}
\noindent
Şimdi \olref[valpha]{defValphas} tanımını doğrulamak için yalnızca
$A=\emptyset$, $\tau(x)=\Pow{x}$ ve $\theta(x)=\bigcup x$ alın. Böylece
nihayet $V_\alpha$'ların tanımı doğrulanmış olur!{}

